\documentclass[11pt]{article}

\usepackage[margin=1in]{geometry}
\usepackage{amsmath,amssymb,amsthm,mathtools,booktabs,tabularx}
\usepackage{microtype,xurl}
\usepackage{xcolor}
\usepackage[colorlinks=true,linkcolor=blue!55!black,
            citecolor=blue!55!black,urlcolor=blue!55!black]{hyperref}

\newtheorem{theorem}{Theorem}[section]
\newtheorem{proposition}[theorem]{Proposition}
\newtheorem{corollary}[theorem]{Corollary}
\theoremstyle{definition}
\newtheorem{definition}[theorem]{Definition}
\theoremstyle{remark}
\newtheorem{remark}[theorem]{Remark}

\newcommand{\dd}{\mathrm d}
\newcommand{\cert}[1]{\nolinkurl{#1}}
\newcolumntype{Y}{>{\raggedright\arraybackslash}X}

\title{What Conformal Gravity Must Assume\\
\large to Explain Galactic Rotation Curves}
\author{Claude Opus 5\thanks{Anthropic model; computations and drafting.}
\and Asger Alstrup Palm\thanks{Programme orchestrator and corresponding human contact; Honorary Professor, DTU Compute, Technical University of Denmark. \texttt{asger@area9.dk}.}}
\date{Working draft, 6 August 2026}

\begin{document}
\maketitle

\begin{abstract}
Conformal (Weyl) gravity admits a static spherically symmetric vacuum
solution whose metric function carries a term linear in \(r\).  That term is
the basis of a long-standing proposal to fit galactic rotation curves without
dark matter.  We ask a reverse-physics question about it: not whether the
fits succeed, but \emph{what must be assumed for them to mean anything}.

Six results, each an exact symbolic computation with machine-checked
controls.  (i) The linear term is not an ansatz: in the conformal gauge the
Bach vacuum equation is exactly the \emph{biharmonic} equation
\(\nabla^{4}B=0\), so \(\gamma r\) is the point-source response of the
governing operator, on the same footing as the Newtonian \(1/r\).  (ii) If
the baryonic Tully--Fisher relation holds, \(\gamma\) is forced to be a
universal constant and is identified with the MOND acceleration scale,
\(\gamma=a_{0}/2c^{2}\); a \(\gamma\) carrying a mass-proportional piece
gives the wrong slope.  (iii) Conformal invariance forces the source to be
trace-free, so a perfect fluid must satisfy \(p=\rho/3\) and pressureless
baryons are excluded outright.  (iv) The standard repair---breaking the
symmetry with a conformally coupled scalar---fails in its simplest form: a
\emph{constant} vacuum expectation value pins the Ricci scalar and thereby
forces \(\gamma=0\), removing the very term at issue.

Two further results say what the surviving branch costs.  (v) \(\gamma\) is
not a conformal invariant: the conformal group acts on the four parameters of
the general solution through a single M\"obius map of the radius, and shifts
\(\gamma\).  The frame in which the scalar is constant is reached by exactly
that map, so for a non-constant scalar the linear term is generated by the
profile, with \(\gamma=-2(S'/S)-3u(S'/S)^{2}\).  The demand that \(\gamma\)
be universal is therefore the demand that the scalar's \emph{fractional
radial gradient} be universal.  (vi) Which frame matter follows is a
variational identity, not an interpretation: a mass generated by the scalar
satisfies \(-\!\int hS\,\dd s_{g}=-\!\int h\,\dd s_{S^{2}g}\), so the mass
scale \emph{is} the conformal factor and such a particle follows geodesics of
the frame where \(\gamma=0\).  A particle whose mass does not track the
scalar follows the other frame, where \(\gamma\) acts.

The conclusion is a trichotomy with one surviving branch, and on that branch a
fork with two computed sides.  We do not adjudicate it.  We locate it, and the
location is now an equation.
\end{abstract}

\section{The question}

Conformal gravity has been proposed as an explanation of galactic rotation
curves without dark matter, on the strength of a term linear in \(r\) in its
static spherically symmetric vacuum solution.  The literature on whether the
fits succeed is substantial and we add nothing to it.

We ask instead the reverse-physics question: \emph{what has to be true for
the explanation to be an explanation?}  A fit with a per-galaxy free
parameter is a description; a fit with a parameter forced by the theory is a
prediction.  Which of these the proposal is depends on facts that can be
settled by computation, and this paper settles four of them.

Every statement below is an exact symbolic identity over \(\mathbb{Q}\),
carried out with unspecified functions rather than sampled fixtures, with
controls designed to fail if the claim were vacuous.  Certificates and
verifier commands are listed in \S\ref{sec:verification}.

\paragraph{Conventions.}
We work in \(D=4\), signature \((-,+,+,+)\), with the static spherically
symmetric ansatz
\[
  \dd s^{2} \;=\; -B(r)\,\dd t^{2} + \frac{\dd r^{2}}{B(r)}
  + r^{2}\dd\Omega^{2},
\]
the conformal gauge \(g_{tt}g_{rr}=-1\).  \S\ref{sec:gauge} addresses what
that gauge costs.  \(B_{ab}\) denotes the Bach tensor.

\section{The linear term is forced, and it is a point-source response}
\label{sec:forced}

\begin{theorem}[Completeness without an ansatz]
\label{thm:complete}
Let \(B\) be any \(C^{4}\) function, nonvanishing on an interval.  Then
\(B_{ab}=0\) identically if and only if
\[
  (rB)'''' = 0
  \qquad\text{and}\qquad
  c_{1}^{2}-3c_{0}c_{2}=1,
\]
where \(B=c_{0}/r+c_{1}+c_{2}r+c_{3}r^{2}\).
\end{theorem}

The proof is a decomposition.  With \(L:=(rB)''''\) and \(N\) the nonlinear
remainder, the three independent Bach rows satisfy, as exact identities in
\(B\) and its derivatives,
\[
  B_{\theta\theta}=\frac{N}{24r^{2}},\qquad
  B_{tt}=\frac{B\,(N+2r^{3}BL)}{24r^{4}},\qquad
  B_{rr}=\frac{-N+2r^{3}BL}{24r^{4}B},
\]
so all three lie in the span of \(\{N,\;r^{3}BL\}\), and both generators are
recoverable from the rows.  Hence \(B_{ab}=0\) forces \(L=0\) and \(N=0\).

The essential point is that \(L\) is \emph{linear}.  Substituting \(u=rB\)
turns \(L=0\) into \(u''''=0\), whose solution space is the cubics by four
integrations---no ansatz.  On that space \(N\) collapses to the constant
\(4(1+3c_{0}c_{2}-c_{1}^{2})\), leaving one algebraic condition.

\begin{theorem}[The vacuum equation is biharmonic]
\label{thm:biharm}
In spherical symmetry \(\nabla^{2}f=f''+2f'/r\), and for unspecified \(B\)
\[
  \nabla^{4}B \;=\; B''''+\frac{4B'''}{r} \;=\; \frac{1}{r}\,(rB)''''.
\]
\end{theorem}

So Theorem~\ref{thm:complete}'s linear equation is \(\nabla^{4}B=0\): the
Bach vacuum condition says \(B\) is \emph{biharmonic}.  The forced solution
space \(\{1/r,\,1,\,r,\,r^{2}\}\) is the radial biharmonic kernel, and
\(r^{3}\) and \(r^{-2}\)---its nearest neighbours---are verified not to lie
in it.

\begin{corollary}[\(\gamma r\) is a point-source response]
For \(\nabla^{2}\) the point-source potential is \(1/r\) and its coefficient
is the mass; this is Newton, and the origin of \(\beta=GM/c^{2}\).  For
\(\nabla^{4}\) the corresponding object is \(r\), since \(\nabla^{2}r=2/r\)
is harmonic away from the origin.
\end{corollary}

This is a stronger statement than ``\(\gamma r\) survives the equation''.  It
is the term the equation \emph{produces} for a localised source, and it sits
on the same footing as the Newtonian term rather than being appended to it.

\begin{remark}
That the general static spherically symmetric vacuum solution of conformal
gravity has this form is classical.  What is contributed here is the removal
of the ansatz and the identification of the operator, both computed for an
unspecified function.
\end{remark}

\section{Tully--Fisher forces \(\gamma\) to be universal}
\label{sec:tf}

With \(\beta=GM/c^{2}\) the locally measured circular velocity is
\(v^{2}/c^{2}=rB'/2B\), which at small field strength is
\(\beta/r+\gamma r/2-kr^{2}\).  In the galactic regime this has a stationary
point, and it is a minimum:
\[
  r_{*}=\sqrt{2\beta/\gamma},
  \qquad
  v^{4}=2GM\gamma c^{2}\Bigl(1-\tfrac{3}{2}\beta\gamma\Bigr).
\]

\begin{theorem}[The conditional]
\label{thm:tf}
To leading order in \(\beta\gamma\): if \(\gamma\) is universal then
\(v^{4}\propto M\), the baryonic Tully--Fisher relation with slope \(4\) in
\(v\).  If \(\gamma\) carries a piece proportional to \(M\) then
\(v^{4}\propto M^{2}\).  Matching the observed \(v^{4}=GMa_{0}\) forces
\[
  \gamma \;=\; \frac{a_{0}}{2c^{2}}.
\]
\end{theorem}

Read in the reverse-physics direction: the observed Tully--Fisher relation,
combined with the forced metric, requires \(\gamma\) to be a universal
constant of nature rather than a per-galaxy parameter, and identifies it with
the MOND acceleration scale.  That is a prediction rather than a fit, and it
is falsifiable in a specific way---a mass-dependent \(\gamma\) moves the
slope.

Two boundaries travel with it.  Identifying \(v_{\min}\) with the observed
flat velocity is an interpretive step, not a theorem.  And the correction
factor \(1-\tfrac{3}{2}\beta\gamma\) is itself mass-dependent, so the slope
is exactly \(4\) only in the limit \(\beta\gamma\to0\); for a galaxy the
departure is of order \(10^{-14}\) and unobservable, but ``exactly linear''
would be too strong.

\begin{remark}
Conformal gravity and MOND are not the same theory and disagree away from
\(r_{*}\): beyond it the conformal curve \emph{rises}
(\(\dd v^{2}/\dd r|_{2r_{*}}=3\gamma/8>0\)) where MOND's are asymptotically
flat.  Only the \(-kr^{2}\) term can bend it back.  The outer rotation curve
is therefore a discriminant, not a detail.
\end{remark}

\section{What may source the theory}
\label{sec:matter}

\S\ref{sec:forced} and \S\ref{sec:tf} concern the vacuum.  \(\gamma\) is a
free coefficient there; the explanatory claim needs it \emph{sourced}.

A conformally invariant action satisfies
\(g_{ab}\,\delta S/\delta g_{ab}=0\).  Applied to the gravitational action
this is the Noether identity of the theory; applied to the matter action it
reads \(T^{\mu}{}_{\mu}=0\).  It is one theorem used twice, and we take it as
given.

\begin{theorem}[Trace obstruction]
\label{thm:trace}
A perfect fluid has \(T^{\mu}{}_{\mu}=3p-\rho\), so it is trace-free if and
only if \(p=\rho/3\).  Dust, \(p=0\), forces \(\rho=0\).
\end{theorem}

Radiation, uniquely, and pressureless matter not at all.  A galaxy's baryons
are pressureless to excellent approximation, so on the conformally invariant
branch a galaxy is not an admissible source.

The standard repair is to break the symmetry with a conformally coupled
scalar acquiring a vacuum expectation value, generating a scale so that
ordinary matter can source the theory after all.  In its simplest form this
fails, and it fails on the term at issue.

\begin{theorem}[A constant vacuum expectation value forces \(\gamma=0\)]
\label{thm:vev}
For the conformally coupled scalar, \(\Box S+\tfrac{1}{6}RS+4\lambda S^{3}=0\)
reduces for constant \(S=S_{0}\neq0\) to \(R=-24\lambda S_{0}^{2}\), a
constant.  For the forced family,
\[
  R \;=\; 12k-\frac{6\gamma}{r}+\frac{2(1-w)}{r^{2}},
\]
which is constant in \(r\) if and only if \(\gamma=0\) and \(w=1\).
\end{theorem}

A constant vacuum expectation value does not merely supply a scale; it pins
the Ricci scalar, and the forced family is compatible with a pinned \(R\)
only where the linear term vanishes.  Note that \(u\), the Newtonian
coefficient, drops out of \(R\) entirely: the obstruction is specific to
\(\gamma\).

\section{The trichotomy}
\label{sec:trichotomy}

Collecting \S\ref{sec:matter}, the possibilities are organised by one
question---\emph{is there a scale, and is it constant or field-dependent?}

\begin{center}
\begin{tabularx}{\textwidth}{@{}lYY@{}}
\toprule
& scale & consequence \\
\midrule
(A) conformal matter & none & source must be trace-free, hence radiation; a galaxy is excluded, so \(\gamma\) is not tied to \(M\) and Tully--Fisher is a coincidence \\
(B1) constant \(\langle S\rangle\) & constant & \(R\) pinned, hence \(\gamma=0\); the linear term is removed \\
(B2) non-constant \(S\) & field-dependent & \(\gamma\) tied to the scalar profile; \S\ref{sec:transfer} makes the tie explicit \\
\bottomrule
\end{tabularx}
\end{center}

(A) does not remove \(\gamma\); it decouples \(\gamma\) from mass, which
destroys the explanation while leaving the term.  (B1) removes \(\gamma\)
outright.  Exactly one branch survives, and it is the one the literature in
fact uses.

\begin{corollary}[Where the assumption lives]
On branch (B2), \(\gamma\) is determined by the scalar profile, which is a
property of each configuration.  Theorem~\ref{thm:tf} requires \(\gamma\) to
be the same constant for every galaxy.  Therefore the requirement transfers
to the profile.
\end{corollary}

That is a question about the symmetry-breaking sector rather than about
gravity.  \S\ref{sec:transfer} carries out the transfer, and it turns out to
land on one derived quantity rather than on the profile as a whole.

\section{What the conformal gauge costs}
\label{sec:gauge}

Sections \ref{sec:forced}--\ref{sec:trichotomy} work in the gauge
\(g_{tt}g_{rr}=-1\).  Two things are needed to justify it.

\begin{proposition}[Conformal covariance]
In \(D=4\), \(B_{ab}[\Omega^{2}g]=\Omega^{-2}B_{ab}[g]\), verified for the
Mannheim--Kazanas family with symbolic parameters and \emph{unspecified}
\(\Omega\).  Hence Bach flatness is a property of the conformal class, and
the classification of \S\ref{sec:forced} classifies classes.
\end{proposition}

\begin{proposition}[Reachability is a first-order ODE]
Setting the new radius to \(\rho=\Omega r\) so the angular part remains
\(\rho^{2}\), the gauge condition becomes
\[
  r\,\Omega' \;=\; \Omega\bigl(\Omega\sqrt{ab}-1\bigr),
\]
first order in \(\Omega\), hence locally solvable wherever \(ab>0\) and
\(r\neq0\).
\end{proposition}

What remains assumed is only global continuation---whether \(\Omega\) stays
positive and finite across the exterior.  That is strictly smaller than the
original assumption and is a question of analysis rather than algebra.

\section{Carrying out the transfer}
\label{sec:transfer}

\S\ref{sec:gauge} and \S\ref{sec:forced} between them force a third statement
that neither makes on its own.  Bach flatness is a property of a conformal
class; the general solution in the gauge \(g_{tt}g_{rr}=-1\) is four
parameters subject to one constraint.  So the conformal group must \emph{act}
on those four parameters.

\begin{theorem}[The conformal action on the forced family]
\label{thm:action}
Staying in the gauge is not a choice but an ODE, \(\dd\hat r/\dd r=(\hat
r/r)^{2}\), whose general solution is the M\"obius map \(\hat
r=r/(1-Cr)\)---one constant.  The induced action is
\[
  u\mapsto u,\qquad
  w\mapsto w+3uC,\qquad
  \gamma\mapsto\gamma-2wC-3uC^{2},\qquad
  k\mapsto k+C\gamma-C^{2}w-C^{3}u,
\]
a one-parameter group preserving \(w^{2}+3u\gamma\) identically.  From a
member with \(\gamma=0\) it produces \(\gamma=-2C-3uC^{2}\).
\end{theorem}

So \(\gamma\) is not an invariant of the Bach vacuum.  It is a coordinate on
the conformal class, and the classification of \S\ref{sec:forced} classifies
classes precisely because of this.

A conformally coupled scalar has weight \(-1\), so the frame in which it is
\emph{constant} is reached by exactly this map with \(\Omega=S/S_{0}\)---and
Theorem~\ref{thm:vev} puts \(\gamma=0\) there.  The transfer follows.

\begin{corollary}[The transfer, in closed form]
\label{cor:transfer}
On branch (B2) the profile realising the frame change is
\(S(r)=S_{0}/(1-Cr)\) with \(C=(S'/S)\) at the origin, and
\[
  \gamma \;=\; -2\,(S'/S)\;-\;3u\,(S'/S)^{2}
  \qquad\text{so that}\qquad \gamma\simeq-2\,(S'/S).
\]
Hence Theorem~\ref{thm:tf}'s demand that \(\gamma\) be one constant for every
galaxy is exactly the demand that \(S'/S\) be one constant for every galaxy.
\end{corollary}

That is more specific than ``is the profile universal?''.  Only one derived
quantity has to be universal, and it has to be so across galaxies whose
masses differ by orders of magnitude, for a field sourced by the
configuration.

\subsection*{Which frame matter follows}

Both frames describe the same Bach vacuum, and only one carries a linear
term.  So the disagreement of \cite{Flanagan2006,Mannheim2006} is not about
the vacuum solution and not about the algebra: it is about which frame matter
follows.  That premise is itself computable.

\begin{theorem}[The mass scale is the conformal factor]
\label{thm:frame}
A point particle of mass \(m\) extremises \(-\!\int m\,\dd s\).  If the mass
is generated by the symmetry-breaking scalar, \(m=hS\), then
\[
  -\!\int hS\,\dd s_{g}
  \;=\;-\!\int h\sqrt{-S^{2}g_{ab}\,\dd x^{a}\dd x^{b}}
  \;=\;-\!\int h\,\dd s_{S^{2}g}
\]
identically, so the trajectory is a geodesic of \(S^{2}g\) and not of \(g\).
Since \(S^{2}g\) is the frame in which \(S\) is constant up to the constant
\(S_{0}^{2}\), such a particle follows geodesics of the frame with
\(\gamma=0\).
\end{theorem}

Both branches are exhibited with their circular-orbit velocities, and both
frames are confirmed Bach vacua, so this is a fork between two admissible
solutions of one theory:

\begin{center}
\begin{tabularx}{\textwidth}{@{}lYY@{}}
\toprule
mass & trajectory follows & linear coefficient of \(v^{2}\) \\
\midrule
\(m\) constant, independent of \(S\) & geodesics of \(g\) & \(\gamma/2\) \\
\(m=hS\), generated by \(S\) & geodesics of \(S^{2}g\) & \(0\) \\
\bottomrule
\end{tabularx}
\end{center}

The first entry reproduces Theorem~\ref{thm:tf}'s weak-field formula by a
different route.  The disagreement is not an artefact of the expansion: the
exact circular-orbit velocities differ too.

A rotation-curve fit needs the first row.  The second row is what the
symmetry breaking was introduced to provide.  Whether a matter sector can
generate a mass scale by breaking conformal invariance while leaving particle
masses independent of the breaking field is a question about that sector, and
we do not answer it.  What is no longer available is leaving the choice
implicit.

\section{What this paper does not establish}

\begin{itemize}\itemsep2pt
\item \textbf{No novelty in the physics.} The Mannheim--Kazanas solution, its
  reduction to a fourth-order Poisson problem, and the trace-free condition
  on conformal matter are all classical.  The contribution is that each is
  computed here with the ansatz removed and the assumptions isolated, rather
  than cited.
\item \textbf{No fits.} No galaxy data enters anywhere.  The only
  observational inputs are the \emph{shape} of the Tully--Fisher relation,
  used as a premise in Theorem~\ref{thm:tf}, and that galaxies are
  pressureless, used in \S\ref{sec:matter}.  Neither is our measurement.
\item \textbf{No verdict on branch (B2).}  \S\ref{sec:transfer} makes the
  branch explicit but does not close it.  Corollary~\ref{cor:transfer} gives
  the profile that \emph{realises the frame change}; it does not show that
  profile solves the coupled scalar-plus-Bach system with matter present,
  which is the remaining half of ``forced versus sourced''.  Nor does it show
  \(S'/S\) cannot be universal---only what universality requires.
\item \textbf{No resolution of the Flanagan--Mannheim dispute}
  \cite{Flanagan2006,Mannheim2006} over the Newtonian limit under conformal
  matter coupling.  Theorem~\ref{thm:frame} narrows it to one premise and
  computes both branches, so each side's commitment is visible; it does not
  say which a real matter sector realises.  The Yukawa form \(m=hS\) is
  itself a premise, standard but entered rather than derived.
\item \textbf{Nothing about the ghost.}  Conformal gravity's fourth-order
  kinetic operator carries an Ostrogradsky ghost.  That is orthogonal to
  everything here and is treated elsewhere in this programme.
\item \textbf{No numerical value.}  \(a_{0}\) is not measured here and no
  published \(\gamma\) is compared; \(\gamma=a_{0}/2c^{2}\) is an algebraic
  consequence of a premise, not a numerical agreement.
\end{itemize}

\section{Verification}
\label{sec:verification}

Each result is a fail-closed certificate with a deterministic verifier.
Controls are designed to fail if the claim were vacuous; several did so
during development and are recorded in the certificates.

\begin{center}
\begin{tabularx}{\textwidth}{@{}lYl@{}}
\toprule
result & certificate & checks \\
\midrule
Thm.~\ref{thm:complete} & \cert{BH0B\_GENERAL\_STATIC\_SPHERICAL\_COMPLETENESS} & 24 \\
Thm.~\ref{thm:tf}       & \cert{BH0C\_TULLY\_FISHER\_SCALING}                   & 14 \\
Thm.~\ref{thm:biharm}   & \cert{BH0D\_BACH\_VACUUM\_IS\_BIHARMONIC}             & 9 \\
Thm.~\ref{thm:trace}    & \cert{BH0E\_MATTER\_TRACE\_DILEMMA}                   & 8 \\
Thm.~\ref{thm:vev}      & \cert{BH0F\_CONSTANT\_VEV\_FORCES\_GAMMA\_ZERO}       & 12 \\
\S\ref{sec:gauge}       & \cert{BH0G\_CONFORMAL\_CLASS\_INVARIANCE}             & 10 \\
Thm.~\ref{thm:action}   & \cert{BH0H\_UNIVERSALITY\_TRANSFER}                   & 18 \\
Thm.~\ref{thm:frame}    & \cert{BH0I\_WHICH\_FRAME\_MATTER\_FOLLOWS}            & 13 \\
\bottomrule
\end{tabularx}
\end{center}

Theorems~\ref{thm:complete}, \ref{thm:vev}, \ref{thm:action} and
\ref{thm:frame}, and \S\ref{sec:gauge}, are computed in \textsc{sympy}; Theorems~\ref{thm:biharm} and \ref{thm:trace} are
computed in Forge over exact rationals with symbolic differentiation of
unspecified functions.  Theorem~\ref{thm:complete}'s conclusion carries an
independent Forge rail using a different arithmetic stack, representation and
code path; that rail confirms which metrics are Bach-flat but not the row
decomposition itself, which remains single-implementation.

\begin{thebibliography}{9}
\bibitem{MK1989} P.~D.~Mannheim and D.~Kazanas,
  \emph{Exact vacuum solution to conformal Weyl gravity and galactic rotation
  curves}, Astrophys.\ J.\ \textbf{342} (1989) 635.
\bibitem{Riegert1984} R.~J.~Riegert,
  \emph{Birkhoff's theorem in conformal gravity}, Phys.\ Rev.\ Lett.\
  \textbf{53} (1984) 315.
\bibitem{Flanagan2006} \'E.~\'E.~Flanagan,
  \emph{Fourth order Weyl gravity}, Phys.\ Rev.\ D \textbf{74} (2006) 023002.
\bibitem{Mannheim2006} P.~D.~Mannheim,
  \emph{Alternatives to dark matter and dark energy}, Prog.\ Part.\ Nucl.\
  Phys.\ \textbf{56} (2006) 340.
\bibitem{Milgrom1983} M.~Milgrom,
  \emph{A modification of the Newtonian dynamics as a possible alternative to
  the hidden mass hypothesis}, Astrophys.\ J.\ \textbf{270} (1983) 365.
\end{thebibliography}

\end{document}
